% ============================================================================
% Teraval — Presentation Script (speaker script for the 7–10 minute talk)
% Compile on Overleaf / locally with pdfLaTeX. No image files needed.
% ============================================================================
\documentclass[12pt,a4paper]{article}

\usepackage[a4paper,top=2.6cm,bottom=2.6cm,left=2.4cm,right=2.4cm]{geometry}
\usepackage{amsmath}
\usepackage{enumitem}
\usepackage{parskip}
\usepackage{xcolor}
\usepackage[hidelinks]{hyperref}

% --- single-line black border on EVERY page (matches main.tex) ---
\usepackage{eso-pic}
\usepackage{tikz}
\usetikzlibrary{calc}
\AddToShipoutPictureBG{%
  \begin{tikzpicture}[remember picture,overlay]
    \draw[line width=0.6pt]
      ($(current page.north west)+(1cm,-1cm)$) rectangle
      ($(current page.south east)+(-1cm,1cm)$);
  \end{tikzpicture}%
}

\definecolor{tvblue}{HTML}{1D4ED8}
\definecolor{tvgreen}{HTML}{047857}
\definecolor{tvamber}{HTML}{B45309}
\definecolor{tvgrey}{HTML}{444444}
\setlength{\parskip}{5pt}
\setlength{\emergencystretch}{2em}  % absorb tiny overruns → no overfull boxes

% block labels
\newcommand{\block}[2]{\par\medskip\noindent\textcolor{#1}{\textbf{#2}}\par\nobreak\smallskip}
\newcommand{\SC}{\block{tvblue}{SCRIPT --- speak this}}
\newcommand{\EX}{\block{tvgreen}{EXPLANATION --- the words used above}}
\newcommand{\PQ}{\block{tvamber}{PROBABLE QUESTIONS \& ANSWERS}}
\newcommand{\qa}[2]{\par\noindent\textbf{Q.}~#1\par\noindent\textbf{A.}~#2\par\smallskip}
\newenvironment{terms}{\begin{description}[leftmargin=0pt,itemsep=2pt,font=\normalfont\bfseries]}{\end{description}}

\begin{document}

% ============================ TITLE PAGE ============================
\begin{titlepage}
\centering
\vspace*{1.4cm}
{\large\scshape Corporate Finance\par}
\vspace{0.4cm}
{\rule{0.6\textwidth}{0.4pt}\par}
\vspace{1.4cm}
{\Huge\bfseries Teraval\par}
\vspace{0.5cm}
{\Large Presentation Script\par}
\vspace{0.3cm}
{\large The words to speak for the 7--10 minute talk\par}
\vspace{1.6cm}
{\large\itshape AI-Enabled Corporate-Finance Capital-Budgeting Dashboard\par}
\vspace{2.0cm}
{\large\textbf{Team}\par}
\vspace{0.3cm}
{\large Kartik Joshi \quad Prem Kukreja \quad Gagandeep Singh\par}
{\large Samuel Alex \quad Aditya Chitale\par}
\vspace{1.4cm}
{\large \textbf{Instructor:} Dr.\ Nathaniel Christopher\par}
\vfill
{\small Submitted for the partial fulfilment of the Corporate Finance module.\par}
\end{titlepage}

% ============================ HOW TO USE ============================
\section*{How to use this script}
This is the \emph{spoken} part of the presentation --- the words to say, not the slides.
It is written so anyone can read it aloud and sound natural, not like they are reading.
Under every section there are three blocks:
\begin{itemize}
  \item \textcolor{tvblue}{\textbf{SCRIPT}} --- the actual words to speak (about 60--90 seconds each).
  \item \textcolor{tvgreen}{\textbf{EXPLANATION}} --- plain-English meanings of any term used, so the
  speaker fully understands what they are saying.
  \item \textcolor{tvamber}{\textbf{PROBABLE QUESTIONS \& ANSWERS}} --- likely questions and short,
  simple answers to have ready.
\end{itemize}
Read in full, the ten SCRIPT blocks run about 11--13 minutes at a calm pace. For a tight 7--10 minute
slot, speak briskly and, if needed, merge Sections~3--4 (concepts + data) and 6--7 (cash flow +
scenarios), or trim the shortest sections --- the ten are modular on purpose. Every number in here
matches the dashboard exactly. Tip: pause after each big number, and point at the screen when you name
a tab.

\tableofcontents
\newpage

% ============================================================================
\section{Opening \& Overview}

\SC
Hi everyone, we're team Teraval, and we built a tool that helps a company make one very expensive
decision using numbers instead of guesswork. Here's the situation. Imagine a company --- we call it
Barq AI --- that wants to build a giant warehouse full of powerful AI computers in Abu Dhabi, and
rent out that computing power. Building it costs about five point eight billion dirhams. That's a
bet-the-company amount of money. So the big question is simple: is it worth it, yes or no? Our
dashboard, on the screen now, answers exactly that. It does all the finance maths automatically, shows
it as clear charts, lets you drag sliders to test ``what if'' situations, layers real artificial
intelligence on top to forecast and stress-test, and ends with a clear recommendation. Over the next
few minutes we'll walk you through it tab by tab. And the short version of our answer, which we'll
prove step by step, is: yes, build it --- but carefully, with a few conditions attached.

\EX
\begin{terms}
  \item[Barq AI:] the (realistic but made-up) company making the decision. ``Barq'' means lightning
  in Arabic --- fitting for a business that sells computing speed.
  \item[Dashboard:] a single screen that shows all the important numbers and charts together.
  \item[AI computers / GPUs:] special chips that do the heavy maths for artificial intelligence; here
  they are rented out by the hour.
\end{terms}

\PQ
\qa{In one line, what does your tool actually do?}{It takes a huge investment decision and turns it
into clear numbers and charts, then recommends yes or no --- with the risks spelled out.}
\qa{Is Barq AI a real company?}{No --- it's a realistic made-up operator, so we can use real market
data without needing a real firm's private financials.}
\qa{Why should a non-finance person care?}{Because it explains every number in plain English and even
has a chatbot you can ask --- you don't need a finance degree to follow the decision.}

% ============================================================================
\section{The Financial Problem}

\SC
So here's the thing we're really deciding. Barq AI needs to spend about five point eight billion
dirhams to build and own a forty-megawatt AI data-centre hall --- that's a very large building full of
computers --- and run it for eight years, renting the computing power to AI companies. Why is this so
hard? Three reasons. First, the money. Five point eight billion is enormous; get it wrong and the
company could sink. Second, it's almost impossible to undo. Once you pour the concrete and install
special cooling, you can't easily turn it back into cash if things go bad. And third --- this is the
scary one --- the price they'd earn is falling fast. The price to rent one GPU for one hour dropped
from about eight dollars to roughly three dollars as lots of new suppliers appeared. If that price
keeps falling, the whole project can lose money. One thing does help, though: Abu Dhabi has very cheap
electricity, which is the biggest running cost, so location is a real advantage.

\EX
\begin{terms}
  \item[40 megawatts (MW):] the amount of electrical power the computers use --- a measure of how big
  the facility is. 40 MW is very large.
  \item[Irreversible:] once you've built it, you can't easily get your money back out.
  \item[GPU-hour:] renting one GPU for one hour --- the unit Barq sells and earns money on.
\end{terms}

\PQ
\qa{Why not just rent computers instead of building?}{Great question --- we actually test exactly that
later; it's a whole tab. Building only wins if the hall stays very busy.}
\qa{Where does the five point eight billion figure come from?}{It's the cost of the GPU hardware plus
the building and cooling --- about three hundred racks of NVIDIA GB200 chips plus a purpose-built
40~MW hall.}
\qa{What's the single biggest risk?}{The GPU rental price falling. It's the number that swings the
answer the most, and it's already close to our danger line.}

% ============================================================================
\section{The Finance Concepts (the words we'll use)}

\SC
Before the numbers, let me quickly explain the finance words, because we'll use them a lot, and
they're simpler than they sound. The first idea is that money today is worth more than money later,
because today's money can be invested --- so we ``discount'' future money to today's value. Add up all
the future cash the project brings in, shrink it to today's value, subtract what we spend, and you get
the N-P-V, the Net Present Value. If it's positive, the project makes us richer --- that's the whole
game. The I-R-R is just the yearly percentage return; if it beats our cost of money, good. The
Profitability Index is the value we get per dirham spent --- above one is good. Payback is how many
years until we get our money back. And the discount rate we use is called the WACC --- basically the
company's yearly cost of money. Everything on the dashboard is built from these few simple ideas.

\EX
\begin{terms}
  \item[Discounting / present value:] shrinking a future amount to what it's worth today, because
  today's money could be earning a return.
  \item[NPV (Net Present Value):] all the future cash, shrunk to today's value, minus what we spend.
  Positive means the project adds value.
  \item[IRR (Internal Rate of Return):] the yearly \% return the project earns on the money tied up
  in it.
  \item[MIRR:] a more careful version of IRR that fixes a known flaw in it.
  \item[PI (Profitability Index):] value created for every one dirham invested; above 1 is good.
  \item[Payback / discounted payback:] how many years to earn the money back --- the discounted one
  also accounts for the time-value of money, so it's longer and more honest.
  \item[WACC:] the blended yearly cost of the money used to fund the project (both borrowing and
  investors' money). It's the bar the project has to beat.
\end{terms}

\PQ
\qa{Why do you discount future money at all?}{Because a dirham next year is worth less than one today
--- today's could already be invested and growing. Ignoring that would overstate the project.}
\qa{Which number is the most important?}{NPV. It directly answers ``does this add value?'' Everything
else supports it.}
\qa{Isn't IRR enough on its own?}{Not quite --- IRR can mislead, so we also show MIRR, and we always
look at NPV alongside it.}

% ============================================================================
\section{Data \& Assumptions}

\SC
Now, where do our numbers come from? We're honest about every input. The key ones are: the GPU rental
price, four dollars an hour; how busy the hall is, eighty percent; the size, forty megawatts; the
up-front cost, about five point eight billion dirhams; the electricity price, fifteen fils a unit;
the cooling efficiency, and the WACC, nine percent, over an eight-year life. And we tag every single
input by where it comes from --- is it historical data, current market data, a forecast, something the
user typed in, or AI-generated? That table is right there in the Ethics and Audit tab. Two of our
inputs are backed by official government data: we pulled the UAE central bank's base rate, three point
six five percent, and the three-month interbank rate, three point nine four percent, straight from the
central bank, and used them to build our discount rate. So this isn't made-up --- it's anchored to
real, current, official numbers.

\EX
\begin{terms}
  \item[Utilization (80\%):] how much of the time the computers are actually rented and earning ---
  our base assumption is 80 percent.
  \item[Fils:] one hundredth of a dirham; so 15 fils is AED 0.15 per unit of electricity.
  \item[PUE (cooling efficiency):] how much extra power the cooling needs; 1.0 is perfect, 1.2 is
  good for this kind of liquid-cooled hall.
  \item[Data tags (H/C/F/U/AI):] Historical, Current, Forecast, User-entered, or AI-generated --- we
  label every assumption so nothing is hidden.
  \item[Central bank base rate / EIBOR:] official UAE interest rates; we used them to build our
  discount rate.
\end{terms}

\PQ
\qa{Why did you choose 9\% for the WACC?}{We built it up from the official central-bank rate of three
point six five percent, added a premium for the equity investors, and blended in the cost of debt.
Nine percent is actually conservative --- a tougher test than the raw maths gives.}
\qa{Is 80\% utilization realistic?}{It's an achievable mix of long-term contracts and on-demand use ---
and we test lower values too. Below about sixty-seven percent the project stops making money.}
\qa{What about the electricity tariff --- is that official?}{We sourced the two interest rates
officially; the exact industrial tariff download was blocked by a portal outage, so we use the
well-known AED 0.15 benchmark and flag it honestly in the data-provenance panel.}

% ============================================================================
\section{Financial Calculations \& Results}

\SC
Okay, the moment of truth --- what do the numbers say? On the KPI cards you'll see the results. The
NPV is positive: about plus one thousand eight hundred and fifty-four million dirhams. In plain terms,
after paying for everything and accounting for time, the project adds roughly one point eight billion
dirhams of value. The yearly return, the IRR, is sixteen point five percent --- well above our nine
percent cost of money. The Profitability Index is one point three one, comfortably above one. It pays
back in five years, six point three if we're strict about it. So on every standard test, it's a yes.
But --- and this is the key finding --- look at the break-even price: three dollars thirty-four per
GPU-hour. That's the price at which the project exactly breaks even. And the market is already at two
eighty-five to three fifty. So we're profitable, but the safety margin is thin. One more thing: the
dashboard even runs a live self-test to prove its own maths is internally consistent.

\EX
\begin{terms}
  \item[KPI cards:] the boxes at the top showing the headline numbers (NPV, IRR, and so on).
  \item[Break-even price (\$3.34):] the GPU rental price at which NPV is exactly zero. Below it, the
  project loses money.
  \item[Self-test:] a live check that the model's own maths obeys the finance rules (for example,
  discounting at the IRR must give an NPV of zero). If it ever broke, a badge turns red.
\end{terms}

\PQ
\qa{Positive NPV and 16.5\% return --- so it's a clear yes?}{On the numbers, yes. The catch is the thin
margin: break-even is three thirty-four and the market's already close to it. That's why our final
answer has conditions.}
\qa{What does the self-test actually prove?}{That our calculator is internally consistent --- it isn't
fudging numbers. It's like showing your working is correct, live, on screen.}
\qa{Payback is 5 years but discounted payback is 6.3 --- why the difference?}{Because the discounted
one also charges for the time-value of money, so it's always longer. Both are within the eight-year
life, which is fine.}

% ============================================================================
\section{Cash Flow \& Build-vs-Rent}

\SC
Let's look at the money flow over time --- this is the Cash Flow tab. Each bar is one year's cash. Year
zero is the big spend, so it's a deep negative. Then money comes in each year. Notice year four dips ---
that's because we replace ageing GPUs halfway through to stay competitive. The line climbing across is
the running total in today's money; where it crosses zero, around year six, that's our payback point.
Now the smart question isn't ``build or do nothing'' --- Barq needs this computing either way. The real
choice is: build and own the hall, or just rent computing from a cloud giant? Because those two options
last different numbers of years, we compare them fairly using something called Equivalent Annual Cost.
And here's the neat insight: building is a big fixed cost, while renting grows with how much you use.
So building only wins if the hall stays very busy --- above about seventy-seven percent. At our base
case it just wins, but only by a thin margin, around plus three hundred and twenty million.

\EX
\begin{terms}
  \item[Free cash flow (the bars):] the actual cash in or out each year.
  \item[Year-4 refresh:] we reinvest to replace older GPUs mid-life, so that year's cash dips.
  \item[Cumulative discounted line:] the running total of cash in today's money; crossing zero is the
  discounted payback.
  \item[EAC (Equivalent Annual Cost):] turns each option's whole-life cost into one fair ``rent-like''
  yearly figure, so options of different lengths can be compared.
  \item[Build-vs-rent crossover (\textasciitilde 77\%):] the busyness level above which owning beats
  renting.
\end{terms}

\PQ
\qa{Why compare against renting instead of doing nothing?}{Because Barq needs the compute regardless.
The honest comparison is against the next-best way of getting it --- that's a textbook ``relevant cash
flow'' test.}
\qa{Why not just always rent, since it's flexible?}{Renting wins when you're not busy, but it's more
expensive per hour when you're very busy. Above about seventy-seven percent usage, owning is cheaper.}
\qa{What's that dip in year four?}{The GPU refresh --- we spend to replace ageing chips halfway through
so the hall stays competitive. It's a real cost, so it's in the NPV.}

% ============================================================================
\section{Scenarios \& Sensitivity}

\SC
No forecast is certain, so we test the decision against good and bad futures. On the Scenarios tab we
show three: an optimistic case where NPV jumps to about plus ten thousand million; our base case at
plus one thousand eight hundred; and a pessimistic case that goes deeply negative, about minus four
thousand million --- a clear reject. So the decision genuinely can flip. Then on the Sensitivity tab,
you can drag sliders for each input and everything updates live. This tornado chart ranks which input
matters most --- and the clear winner is the GPU rental price, with utilization second. Everything
hinges on those two. We also run a Monte-Carlo simulation --- basically five thousand random ``what-if''
worlds --- and it tells us there's about a twenty-one percent chance the project loses money. So it's a
yes, but a yes with real, measurable risk. That honesty is the point: we're not hiding the downside,
we're quantifying it.

\EX
\begin{terms}
  \item[Scenario:] a deliberate ``what if'' picture --- optimistic, base, or pessimistic.
  \item[Sensitivity / sliders:] dragging an input to see how the answer changes.
  \item[Tornado chart:] a bar chart that ranks the inputs by how much each one moves the NPV.
  \item[Monte-Carlo simulation:] running thousands of random combinations to see the range of possible
  outcomes and the chance of a loss.
  \item[P(loss) \textasciitilde 21\%:] the probability the NPV comes out negative.
\end{terms}

\PQ
\qa{What if GPU prices keep falling?}{That's exactly the risk. Our break-even is three thirty-four, and
the market is already at two eighty-five to three fifty. Below break-even, NPV turns negative --- which
is why our number-one condition is to lock in a rate above three thirty-four before spending.}
\qa{A twenty-one percent chance of loss --- isn't that high?}{It's meaningful, which is why we don't
call it an unconditional yes. Our conditions are designed to shrink exactly that risk.}
\qa{Which input should management watch most?}{The GPU rental price, then utilization --- the tornado
chart shows they dominate everything else.}

% ============================================================================
\section{AI \& Forecasting}

\SC
Now the AI part --- and this is real AI, not just a label. We have seven AI features, but let me
highlight the two trained machine-learning models, because these actually learn from data. The first
is a forecaster: it studies the pattern of GPU prices and predicts where they're heading --- and it's
accurate, with a test error of only about six cents. It forecasts prices settling near two dollars
seventy-three, which is below our break-even --- a genuine warning. The second is a risk classifier
that we trained on thousands of examples to predict ``value-creating or not,'' and it's about
ninety-seven percent accurate on data it has never seen, with a near-perfect score. Importantly, both
models are trained and tested on a proper split of data, and both agree with our finance engine ---
they augment it, they don't replace it. On top of those, we have Monte-Carlo simulation and a grounded
chatbot you can ask questions to. And crucially, the AI never invents numbers --- it's handed the real
computed figures and only explains them.

\EX
\begin{terms}
  \item[Machine-learning model:] a program that learns patterns from example data, rather than being
  told fixed rules.
  \item[Train/test split:] we teach the model on part of the data and check it on a separate part it
  has never seen, to prove it really learned.
  \item[Forecaster (test error \textasciitilde 6\%):] predicts future GPU prices; its error on unseen
  data is tiny.
  \item[Risk classifier (\textasciitilde 97\% accurate):] predicts whether a set of inputs will create
  value; trained on thousands of engine-labelled examples.
  \item[Grounded chatbot:] an AI assistant that can only use the real numbers it's given --- so it can't
  make things up.
\end{terms}

\PQ
\qa{Is this really machine learning, or just charts?}{Really machine learning --- two models that train
on data, are tested on a held-out set, and hit a stated accuracy bar. It's reproducible: same data,
same result every time.}
\qa{If the finance engine already gives the answer, why add ML?}{The engine gives the exact answer for
given inputs; the ML forecasts the uncertain inputs (like future price) and gives a fast risk read.
It augments the engine, it doesn't replace it.}
\qa{How do you stop the chatbot from making things up?}{We hand it the real computed numbers and tell
it to only explain them --- and we added guardrails so it politely refuses anything off-topic.}

% ============================================================================
\section{Board Review \& Ethics}

\SC
Big decisions aren't made by one person, so we built a Board Review tab that mimics a real committee.
Five departments --- Finance, Operations, Commercial, Sustainability, and Technology --- each score the
project out of a hundred using a fixed, transparent rule, and take a stance. Then a weighted verdict
combines them. At our base case, the board says: approve --- with conditions --- scoring sixty-five out
of a hundred. And every score moves live as you change the sliders. On ethics, we take it seriously.
We show forecasts as ranges, never as false certainty. We ground the chatbot so it can't hallucinate
numbers. We stress the downside honestly with the pessimistic case and Monte-Carlo. We tag every
assumption with its source, and our interest rates come straight from the central bank. And most
importantly, we're clear on responsibility: the AI advises, but the final invest-or-reject decision
always stays with the human --- the CFO and the committee, not the machine.

\EX
\begin{terms}
  \item[Board Review:] a view where five departments each score the project and a weighted verdict
  combines them --- like a real approval committee.
  \item[Approve with conditions:] a yes, but only if certain requirements are met first.
  \item[Hallucination:] when an AI confidently makes up a wrong fact --- we prevent it by grounding.
  \item[Data provenance:] a record of exactly where each official number came from.
\end{terms}

\PQ
\qa{Why ``approve with conditions'' and not just approve?}{Because several departments attach a
non-negotiable requirement --- like locking in customers first. Any unmet requirement means it's a
conditional yes, never a blank cheque.}
\qa{Who is responsible if the decision goes wrong?}{The human decision-makers --- the CFO and the
investment committee. Our tool supports the decision; it never makes it.}
\qa{How do you handle bias in the AI forecasts?}{We stress-test the downside deliberately with the
pessimistic scenario and a five-thousand-run simulation, and we lead with the probability of a loss.}

% ============================================================================
\section{Final Recommendation \& Closing}

\SC
So, to bring it all together --- what should Barq AI do? Our recommendation is: accept the project, but
conditionally. It creates about one point eight billion dirhams of value with a healthy sixteen point
five percent return, so the upside is real. But the margin is thin and there's a one-in-five chance of
a loss, driven almost entirely by the GPU rental price. So we attach three conditions. One: lock in
customers with multi-year contracts at a rate above our three-dollar-thirty-four break-even before
spending. Two: release the money in stages, tied to the hall actually staying busy, instead of all at
once. And three: keep a ``rent instead'' fallback ready in case the market turns. Do those three
things, and a thin, risky yes becomes a smart, protected yes. That's Teraval --- we turned a five point
eight billion dirham gamble into a clear, defensible, data-backed decision. Thank you --- we're happy
to take any questions.

\EX
\begin{terms}
  \item[Conditional accept:] yes, provided specific protections are put in place first.
  \item[Contracted offtake:] signing customers to multi-year deals that guarantee a price and usage.
  \item[Staged capex:] releasing the investment money in steps, each unlocked by hitting a target,
  rather than all up front.
\end{terms}

\PQ
\qa{If you had to give a one-word answer, yes or no?}{Yes --- but a conditional yes. The conditions are
what turn a risky bet into a sensible investment.}
\qa{What's the single most important condition?}{Locking in customer contracts above the three
thirty-four break-even before committing the money --- it directly removes the biggest risk.}
\qa{How confident are you in the recommendation?}{Confident, because it's not opinion --- it's built on
verified maths, real official data, stress-testing, and trained models, all of which point the same
way. And we've been honest about the one-in-five downside.}

\vspace{1em}
\begin{center}\small\textcolor{tvgrey}{%
End of script. All ten SCRIPT blocks read in full run about 11--13 minutes; trim or merge sections to
fit a 7--10 minute slot.}\end{center}

\end{document}
